% Page 040

\seite{40}

\abschnitt{Revision}
\pstart
Am 27.\ Oktober wurde die hiesige Schule durch\ob
Herrn Regierungs- und Schulrat Bettermann\unsicher\ob
revidiert.

\pend
\abschnitt{Vertretung}
\pstart

Für den Monat Dezember erhielt Herr Lehrer Fritz\ob
aus Seffern wegen Krankheit Urlaub. Der hiesige\ob
Lehrer mußte abwechselnd während dieser Zeit den\ob
Unterricht in Seffern und Sefferweich abhalten.

\pend
\abschnitt{Schulfeierlichkeiten}
\pstart

Der Geburtstag Seiner Majestät des Kaisers wurde\ob
in der Schule von Seffern gemeinschaftlich von den\ob
4 Schulen der Pfarrei unter Beisein des Herrn\ob
Ortsschulinspektors gefeiert. Während dem Vortrag\ob
einiger Gedichte wurden gemeinschaftlich von den\ob
Schulkindern Lieder gesungen. Der Herr Lehrer\ob
Klein\unsicher{} mahnte zur ergebensten Anhänglichkeit an\ob
Kaiser und Vaterland. Vor Austheilung der Kaiser\ohb
wecken hielt Herr Ortsschulinspektor Pfarrer\ob
Schmitt eine Ansprache etc.

\pend
\abschnitt{Ges. 18/1. 05}
\pstart

Der Ortsschulinspector
Schmitt, Pfr.

\pend
\abschnitt{Osterprüfung}
\pstart

Am 28.\ März wurde die Osterprüfung durch den\ob
Herrn Ortsschulinspector Pfarrer Schmitt aus\ob
Seffern abgehalten.

\pend
\jahresueberschrift{Das Schuljahr 1905/06}
\pstart

Entlassen wurden zu Ostern 1905: 2 Knaben,\ob
von denen einer auf das Gymnasium nach Prüm\ob
geht, der andere wird dem Berufe seines Vaters\ob
dienen. 3 Mädchen wurden entlassen, die sämtlich\ob
zu Hause bei der Haus- und Feldarbeit thätig\ob
bleiben. Die Schülerzahl beträgt nunmehr\ob
29 Knaben u.\ 31 Mädchen. Im Ganzen sind es\ob
60 Kinder. Aufzunehmen sind 9 Knaben u.\ob
7 Mädchen. Alle sind katholisch und vollsinnig\unsicher.

\pend
\abschnitt{Revision}
\pstart

Am 14.\ April wurde die hiesige Schule durch\ob
den Herrn Kreisarzt von Bitburg mit Bezug auf\ob
Reinlichkeit und Ausstattung der Schulräume,\ob
sowie auf Reinlichkeit und Gesundheitszustand\ob
der Kinder revidiert.
\pend
